\chapter{Introduction}
\label{chap:introduction}

\section{Context and Motivation}
In finance, one domain naturally sits at the bleeding edge of software
performance: high-frequency trading (HFT), where firms compete to execute
orders at the lowest possible latency~\citep{aldridge2013hft}. In this context,
the order book is one of the most important components, because it maintains
the state of buy and sell orders for each instrument and sits on the critical
processing path~\citep{gould2013lob}. The design of the order book can therefore
directly affect execution quality and overall system behaviour.

Given this dependence on low-latency software, understanding the performance
implications of different order-book designs is important for both researchers
and practitioners. This dissertation evaluates and compares order-book data
structures in a realistic benchmarking environment. Detailed technical
definitions are provided in Chapter~\ref{chap:background}; this chapter focuses
on the research problem, objectives, and scope of the study.

\section{Problem Statement and Research Gap}
Although there has been significant work on order-book data structures, much of
industry practice remains opaque because implementation details are commercially
sensitive. Any advantage gained from an optimised order-book design can be worth
substantial financial value, so firms rarely disclose architecture decisions in
full detail. This creates a useful role for academic benchmarking: transparent,
reproducible evaluation under controlled conditions.

Another gap is methodological. Many comparisons emphasise isolated algorithmic
performance but do not capture full end-to-end latency, where transport and
system overhead can dominate. This dissertation addresses that gap with a
two-mode approach: direct mode to isolate order-book behaviour, and gateway mode
to evaluate the complete software path from network ingress to matching-engine
processing.

\section{Aim and Objectives}
This dissertation aims to design, implement, and evaluate a C++20 exchange 
infrastructure for benchmarking order book data structures in high-frequency 
trading environments. The specific objectives are as follows:

\begin{itemize}
  \item Define a modular benchmarking framework with interchangeable order-book
  implementations under a common interface.
  \item Implement and validate five order-book variants (map, array, vector,
  pool, and hybrid) with semantic parity across shared test scenarios.
  \item Evaluate performance in direct mode and gateway mode to separate
  algorithmic behaviour from end-to-end system overhead.
  \item Decompose gateway latency into network latency, queue latency, and
  engine latency to localise bottlenecks.
  \item Interpret observed differences using data-structure mechanics and
  hardware-aware reasoning grounded in measured evidence.
\end{itemize}

To evaluate these objectives, the study uses direct mode and gateway mode, and
reports mean latency, P99 latency, throughput, and component-level decomposition
of network latency, queue latency, and engine latency against
\ref{hyp:h1}--\ref{hyp:h3}.


\section{Research Hypotheses}
The evaluation tests three hypotheses that connect algorithmic structure,
allocator strategy, and systems overhead:

\begin{itemize}
  \item \textbf{\ref{hyp:h1} (Data Structure and Cache Locality):} index-based
  order-book layouts are expected to outperform pointer-heavy alternatives in
  isolated operation.
  \item \textbf{\ref{hyp:h2} (Memory Allocation Limits):} memory pooling is
  expected to reduce selected allocation costs but not remove the fundamental
  structural penalties of slower layouts.
  \item \textbf{\ref{hyp:h3} (System-Level Masking):} in gateway operation,
  network and queue overheads are expected to compress or hide many internal
  data-structure performance differences.
\end{itemize}

For clarity, \ref{hyp:h1} is treated as supported when index-based designs
consistently outperform pointer-heavy designs in direct mode across the core
scenarios; \ref{hyp:h2} is treated as supported when pooling improves selected
allocation-heavy cases without eliminating structural penalties across mixed
workloads; and \ref{hyp:h3} is treated as supported when gateway decomposition
shows network and queue costs dominating engine cost, with compressed
cross-implementation differences.

These hypotheses and acceptance criteria are defined formally in
Chapter~\ref{chap:testing_evaluation}.


\section{Scope and Boundaries}
This dissertation focuses on software-path benchmarking within a controlled host
environment, including direct and gateway execution modes, multiple workload
profiles, and comparative analysis across five order-book implementations. The
scope includes functional validation, throughput and latency measurement, and
component-level decomposition of gateway costs. The scope does not include
full physical network-path modelling (for example switch traversal and NIC
hardware contention), and absolute latency values should therefore be interpreted
as environment-bounded rather than universally transferable.


\section{Contributions Overview}
The dissertation makes four contributions. First, it presents a modular,
extensible C++20 exchange framework with interchangeable order-book
implementations and a TCP gateway path for end-to-end benchmarking. Second, it
provides a reproducible benchmarking suite that separates direct-mode and
gateway-mode behaviour while decomposing latency and throughput across pipeline
components. Third, it reports empirical evidence on the trade-offs between
cache-friendly indexing and pointer-based alternatives, including the bounded
benefit profile of memory pooling. Fourth, it demonstrates a systems-level
result: in gateway operation, network and queue overheads can dominate total
latency and mask differences that are visible in isolated mode.


\section{Dissertation Structure}
Chapter~\ref{chap:background} reviews the technical background required for the
study, including market microstructure and hardware-aware performance
considerations. Chapter~\ref{chap:system_design} presents the architectural
design, requirements framework, and benchmarking philosophy. Chapter~\ref{chap:implementation}
details implementation choices for the exchange pipeline and order-book
variants. Chapter~\ref{chap:testing_evaluation} reports functional validation
and performance results against the research hypotheses. Finally,
Chapter~\ref{chap:conclusions} summarises findings, limitations, and future
work.
